\section{第42届复赛第4题：平面激光场中的相对论电子加速}
\subsection{题干}
题目把强激光脉冲简化为沿 $z$ 轴传播、沿 $x$ 轴线偏振的单色平面电磁波：
\[\boldsymbol E=E_0\cos(kz-\omega t)\boldsymbol e_x,\qquad
\boldsymbol B=\frac{E_0}{c}\cos(kz-\omega t)\boldsymbol e_y.\]
电子初始静止时，要求由相对论动力学求动量、能量和位置的隐函数关系；随后讨论给定初始能量、沿 $x$--$z$ 平面入射时的最大增能；最后按给定强度和波长估计辐射阻尼造成的一个周期相对能量损失。
\begin{figure}[htbp]\centering\begin{tikzpicture}[>=Stealth,scale=.9]
  \draw[->,thick] (-3.8,0)--(4.2,0) node[right] {$z$};
  \draw[->,thick] (-3.5,-1.5)--(-3.5,1.7) node[above] {$x$};
  \foreach \z in {-2.8,-1.4,0,1.4,2.8} {
    \draw[blue,thick,->] (\z,-.65)--(\z,.65);
    \node[red] at (\z+.28,.38) {$\odot$};
  }
  \draw[very thick,->] (-3.4,0) .. controls (-2.4,.9) and (-.8,-.8) .. (1.2,.2)
    .. controls (2.2,.7) and (3.1,.35) .. (3.8,.15);
  \fill (-3.4,0) circle (1.5pt) node[below left] {电子};
  \node[blue] at (-2.7,-1.05) {$\boldsymbol E\parallel\boldsymbol e_x$};
  \node[red] at (2.1,1.05) {$\boldsymbol B\parallel\boldsymbol e_y$（出纸面）};
  \node[below] at (.2,-1.5) {解析辅助图：平面波场方向与电子在 $x$--$z$ 平面的运动};
\end{tikzpicture}
\caption{第42届复赛第4题的解析辅助图：题面未给独立题图，本图仅示意场方向和运动平面。}\end{figure}
\subsection{参考解答整理}
洛伦兹力方程为
\[\frac{\mathrm d\boldsymbol p}{\mathrm dt}=-e(\boldsymbol E+\boldsymbol v\times\boldsymbol B),\qquad
\epsilon=\sqrt{m_e^2c^4+p^2c^2}.\]
由 $kc=\omega$ 和初始静止条件，答案给出横向动量
\[p_x=\frac{eE_0}{\omega}\sin(kz-\omega t),\qquad p_y=0.\]
能量与纵向动量满足守恒关系 $\epsilon-cp_z=m_ec^2$，于是
\[p_z=\frac{p_x^2}{2m_ec},\qquad
\epsilon=m_ec^2+\frac{e^2E_0^2\sin^2(kz-\omega t)}{2m_e\omega^2}.\]
令 $\tau=z-ct$，电子位置可写成隐函数；其中
\[x=\frac{eE_0[\cos(kz-\omega t)-1]}{m_e\omega^2},\qquad y=0,\]
并由 $\int_0^\tau\sin^2(k\tau)\,\mathrm d\tau=\tau/2-\sin(2k\tau)/(4k)$ 得到 $z$ 的相应关系。第（2）问在固定初始能量下，以 $\theta_0$ 参数化初始动量并对相位取极值；完整长式见答案第12至13页，本整理不转抄未经逐项排版复核的长分式。

第（3）问先由能流密度的周期平均得到
\[I=\left\langle|\boldsymbol S|\right\rangle=\frac{E_0^2}{2\mu_0c},\qquad E_0^2=2\mu_0cI.\]
代入辐射阻尼率并积分一个周期，答案写为
\[\frac{\Delta\epsilon}{\epsilon_0}=-\frac{1}{1+\dfrac{2\pi m_e^4c^5}{e^4\mu_0^2I\epsilon_0\lambda}}\approx-40\%.\]
\subsection{转写依据}
题干依据 \texttt{CPHO42-Q04\_transcription\_v1.md}；公式与分值依据答案第10至14页和 \texttt{CPHO42-Q04\_formula\_check\_v1.csv}。题面无独立图件，插图为明确标注的解析辅助图，不嵌入原始页面截图。
